\section{Discussion}
\label{sec:discussion}

\subsection{Capability vs.\ Default Behavior}

Our central finding---that \gptfull achieves 0.027 entropy with \basic prompting but 0.977 with \chainofthought on the same sets---has implications beyond this specific task.
It demonstrates that \textbf{capability does not imply default behavior}.
The knowledge required to split a set of fruits into two halves (properties like name length, color, growing region) is clearly present in the model, but it is not activated by a generic question-generation prompt.
This aligns with findings from QuestBench \citep{li2025questbench}, where models can solve fully-specified problems but fail to ask the right clarifying question.

The gap suggests that RLHF or instruction tuning may have biased \gptfull toward generating ``interesting'' or ``relevant'' questions rather than maximally informative ones.
A question like ``Is this a fruit?'' is perfectly reasonable in many contexts---it is just useless for distinguishing among fruits.

\subsection{When Chain-of-Thought Helps and Hurts}

\chainofthought prompting is the most reliable strategy for 8--16 item sets, achieving 100\% perfect splits on structured sets and 65--70\% on homogeneous ones.
The reasoning step forces the model to enumerate candidate properties, count matches, and verify the split before committing to a question.
This explicit verification acts as a self-correction mechanism.

However, \chainofthought degrades on large sets (29 items in \bigbench).
The model's counting becomes unreliable over longer lists, producing incorrect split predictions and suboptimal questions.
For large sets, \explicit prompting is preferable---it communicates the objective without requiring the model to track membership across many items.
This suggests a practical guideline: use \chainofthought for sets of up to $\sim$16 items and \explicit for larger sets.

\subsection{What Makes a Set Hard to Split?}

Our difficulty analysis identifies three factors:

\para{Homogeneity.}
The strongest predictor of difficulty is whether items share a category.
Mixed-category sets are trivially split by category membership.
Single-category sets require the model to find subtler properties---a fundamentally harder task because the space of possible splitting properties is larger and less salient.

\para{Feature availability.}
Among homogeneous sets, difficulty correlates with the availability of salient binary properties.
Animals can be split by habitat (land vs.\ water), diet (herbivore vs.\ carnivore), or size (large vs.\ small).
Fruits have fewer such properties, making them harder to split.
The model sometimes resorts to properties like name length (``Does the name have more than six letters?''), which are creative but fragile---small changes in the item set can break the split.

\para{Set size.}
Larger sets are harder for \chainofthought (which requires tracking all items) but easier for generic questions (which are more likely to produce even splits by chance).
This creates a tradeoff where the optimal strategy depends on set size.

\subsection{Limitations}
\label{sec:limitations}

\para{Judge model bias.}
Using \gptfull as both generator and judge creates a potential consistency bias.
The generator may produce questions that it knows \gptfull-as-judge will answer in a balanced way, inflating entropy scores relative to what a different judge would produce.
A human evaluation or cross-model judging setup would strengthen the results.

\para{Single-run evaluation.}
Each trial was run once with temperature 0.3 for generation.
While the low temperature and fixed seed provide reproducibility, multiple runs with varied temperatures would reveal the variance of question quality.

\para{Limited model coverage.}
We evaluate only OpenAI models.
Testing Claude, Gemini, and open-source models would provide a broader picture of the capability landscape.

\para{Annotation ambiguity.}
Some generated questions have genuinely ambiguous answers (\eg ``Is a tomato a fruit?''---botanically yes, culinarily no).
Our single-judge setup forces binary answers but may mask underlying uncertainty.

\para{Single-question scope.}
We evaluate only the first question in a potential dialogue.
Real-world information seeking requires maintaining balanced splits across multiple rounds, which involves additional challenges like updating the candidate space.
